%----------------------------------------------------------------------------------------
%    PACKAGES AND THEMES
%----------------------------------------------------------------------------------------

\documentclass[aspectratio=169,xcolor=dvipsnames]{beamer}

\AtBeginSection[]
{
    \begin{frame}
        \frametitle{Visão Geral}
        \tableofcontents[currentsection]
    \end{frame}
}

\usefonttheme{professionalfonts} % using non standard fonts for beamer
\usefonttheme{serif} % default family is serif
\usepackage{fontspec}
\setmainfont{XCharter}[
    UprightFont    = *-Roman,
    BoldFont       = *-Bold,
    ItalicFont     = *-Italic,
    BoldItalicFont = *-BoldItalic,
]

\usepackage{hyperref}
\usepackage{graphicx} % Allows including images
\usepackage{booktabs} % Allows the use of \toprule, \midrule and \bottomrule in tables
\input{./slides/SimplePlusTheme/beamerthemeSimplePlus.sty}
\input{./slides/SimplePlusTheme/beamerinnerthemeSimplePlus.sty}
\input{./slides/SimplePlusTheme/beamerfontthemeSimplePlus.sty}
\input{./slides/SimplePlusTheme/beamercolorthemeSimplePlus.sty}

\usepackage{listings}
\setbeamercovered{transparent}

\usepackage{biblatex}
\addbibresource{slides/reference.bib}
\renewcommand*{\bibfont}{\footnotesize}
% Tamanho dos slides de "Leituras Recomendadas": maior que a bibliografia
% completa (\bibfont acima, que rende \tiny via a redefinição de \footnotesize)
% para destacar as leituras curadas.
\newcommand{\leiturasize}{\small}

\usepackage[brazilian ]{babel}
\usepackage{csquotes}
\usepackage{pgfplots}
\pgfplotsset{compat=1.18}

\usepackage{emoji}
\usepackage{amsmath}
\usepackage[xcharter,bigdelims,vvarbb]{newtxmath}
% newtxmath's operators font requests the fontspec family `XCharter(0)' in OT1;
% without these substitutions math digits and operator names silently fall back
% to Computer Modern (LaTeX Font Warning: Font shape `OT1/XCharter(0)/m/n' undefined).
\DeclareFontFamily{OT1}{XCharter(0)}{}
\DeclareFontShape{OT1}{XCharter(0)}{m}{n}{<->ssub * XCharter-TLF/m/n}{}
\DeclareFontShape{OT1}{XCharter(0)}{m}{it}{<->ssub * XCharter-TLF/m/it}{}
\DeclareFontShape{OT1}{XCharter(0)}{b}{n}{<->ssub * XCharter-TLF/b/n}{}
\DeclareFontShape{OT1}{XCharter(0)}{bx}{n}{<->ssub * XCharter-TLF/b/n}{}

\usepackage{bm}
\usepackage[table]{xcolor}
\usepackage{xcolor}
\usepackage{algpseudocode}

\usepackage{cancel}
\usepackage{ulem}

\usetikzlibrary{shapes}
\usetikzlibrary{arrows}
\usetikzlibrary {arrows.meta}
\usetikzlibrary{calc,positioning}
\usepgfplotslibrary{fillbetween}
\usetikzlibrary{angles,quotes}
\usetikzlibrary{tikzmark}


\definecolor{PastelRed}{HTML}{FFC1CC}
\definecolor{PastelBlue}{HTML}{C2F0FF}
\definecolor{PastelBlue2}{HTML}{3182BD}

\definecolor{PastelPink}{HTML}{FFCBE1}
\definecolor{PastelGreen}{HTML}{D6E5BD}
\definecolor{PastelYellow}{HTML}{F9E1A8}
\definecolor{PastelBlue}{HTML}{BCD8EC}
\definecolor{PastelPurple}{HTML}{DCCCEC}
\definecolor{PastelOrange}{HTML}{FFDAB4}


\renewcommand{\footnotesize}{\tiny}

\newcommand{\mespor}{
  \ifcase\month
    \or janeiro
    \or fevereiro
    \or março
    \or abril
    \or maio
    \or junho
    \or julho
    \or agosto
    \or setembro
    \or outubro
    \or novembro
    \or dezembro
  \fi
}
% \usepackage{csquotes}
% \usepackage{minted}
% \usemintedstyle{xcode}
\definecolor{vscLightBackground}{HTML}{FFFFFF}
\definecolor{vscLightText}{HTML}{000000}
\definecolor{vscLightGreen}{HTML}{008000}        % For Comments AND Numbers
\definecolor{vscLightRed}{HTML}{A31515}          % For Strings
\definecolor{vscLightTeal}{HTML}{267F99}
\definecolor{vscBlue}{HTML}{0000FF}% For main keywords (if, for, def...)
\definecolor{vscLightPytorch}{HTML}{800080}      % For PyTorch keywords (dark magenta)

\lstdefinestyle{vscode-light}{
    language=Python,
    backgroundcolor=\color{vscLightBackground},
    basicstyle=\ttfamily\scriptsize\color{vscLightText},
    keywordstyle=\color{vscLightTeal},
    commentstyle=\color{vscLightGreen},
    stringstyle=\color{vscLightRed},
    % Style for PyTorch keywords (using keyword group 2)
    keywordstyle=[2]\color{vscLightPytorch}, 
    morekeywords={self, True, False, None},
    morekeywords=[2]{
        torch, Tensor, device, dtype, to, nn, optim, autograd, utils, data,
        Module, Linear, Conv2d, ReLU, CrossEntropyLoss, Adam, SGD,
        Dataset, DataLoader, randn, zeros, ones, arange, cat, stack,
        view, reshape, sum, mean, max, min, matmul, no_grad, backward
    },
    keywordstyle=[3]\color{vscBlue}, 
    morekeywords=[3]{
        model, loss, y, y_hat, x , optimizer
    },
    frame=single,
    tabsize=1,
    showstringspaces=false,
    breaklines=true,
    captionpos=b,
    extendedchars=true,
    numbers=left,
}

\usepackage[cache=true]{minted}
% minted v3+ (TeX Live 2024+) auto-detects the output directory.
% minted v2 (Ubuntu apt TeX Live 2023) needs it set explicitly to
% match latexmkrc's $out_dir = 'build'.
\makeatletter
\@ifpackagelater{minted}{2024/01/01}{}{%
  \def\minted@outputdir{build/}%
}
\makeatother
\setminted{
    fontsize=\scriptsize,
    style=vs,
    breaklines=true,
    numbers=left,
    numbersep=5pt,
    frame=single,
}
\providecommand{\citeauthoryear}[1]{(\citeauthor{#1},~\citeyear{#1})}

\providecommand{\thetav}{\bm{\theta}}

% vetor coluna 2D com colchetes
\providecommand{\cvec}[2]{\begin{bmatrix} #1 \\ #2 \end{bmatrix}}

\providecommand{\varvector}[1]{\mathbf{#1}}
% instance
\providecommand{\X}{\varvector{X}}
\providecommand{\mlinstance}{\varvector{X}}
\providecommand{\mlinstancei}{\varvector{X}^{(i)}}
\providecommand{\mlinstancex}[1]{\varvector{X}^{(#1)}}

% target
\providecommand{\y}{\varvector{y}}
\providecommand{\mltarget}{\varvector{y}}
\providecommand{\mltargeti}{\varvector{y}^{(i)}}
\providecommand{\mltargetx}[1]{\varvector{y}^{(#1)}}

%hats
\providecommand{\yhat}{\hat{\varvector{y}}}
\providecommand{\mltargethat}{\hat{\varvector{y}}}
\providecommand{\mltargethati}{\hat{\varvector{y}}^{(i)}}
\providecommand{\mltargethatx}[1]{\hat{\varvector{y}}^{(#1)}}

% linear reg
\providecommand{\mllinearregression}{\mltargethat=\mlinstance\bm{\theta}}

\providecommand{\tikzarrowcolor}[3]{%
  \begin{tikzpicture}[remember picture, overlay]
    % Arrow body
    \path[fill=#3] ([shift={(#1,#2)}]current page.south west) 
      rectangle ++(0.3, -0.2);
    % Arrow head
    \path[fill=#3] ([shift={(#1+0.3,#2+0.1)}]current page.south west) -- 
         ++(0, -0.4) -- 
         ++(0.2, 0.2) -- 
         cycle;
  \end{tikzpicture}%
}

\providecommand{\tikzarrow}[2]{%
    \tikzarrowcolor{#1}{#2}{Red}
}



\DeclareMathOperator*{\argmax}{argmax}
\DeclareMathOperator*{\argmin}{argmin}


\usetikzlibrary{spy}

%----------------------------------------------------------------------------------------
%    TITLE PAGE
%----------------------------------------------------------------------------------------

\title{Deep Learning}
\subtitle{Introdução a Redes Convolucionais}

\author{Lucas Silveira Kupssinskü}

\institute
{
    Machine Learning Theory and Applications Laboratory \\
    PUCRS
}
\date{\mespor de \the\year}

%----------------------------------------------------------------------------------------
%    PRESENTATION SLIDES
%----------------------------------------------------------------------------------------

\begin{document}

\begin{frame}
    \titlepage
\end{frame}

\begin{frame}{Overview}
    \tableofcontents
\end{frame}

%========================================================================
\section{Motivação}
%========================================================================

\begin{frame}{Motivação Histórica}
\begin{columns}
    \begin{column}{.6\linewidth}
        \begin{itemize}
            \item David Hubel e Torsten Wiesel realizaram experimentos implantando eletrodos no córtex visual de gatos\footfullcite{hubel1962receptive}.
            \item Aparentemente o processamento começa identificando padrões simples, como arestas.
            \begin{itemize}
                \item \href{https://www.youtube.com/watch?v=jw6nBWo21Zk}{Vídeo}
            \end{itemize}
            \item Camadas seguintes combinam esses padrões em estruturas mais complexas.
        \end{itemize}
    \end{column}
    \begin{column}{.4\linewidth}
        \centering
        \includegraphics[width=\linewidth]{slides/img/02-cnn/hubel_wiesel.jpeg}
    \end{column}
\end{columns}
\end{frame}

\begin{frame}{Redes Convolucionais}
\begin{itemize}
    \item São redes que incluem operações de convolução.
    \begin{itemize}
        \item Usadas principalmente em imagens.
    \end{itemize}
    \item Criadas em 1989 por Yann LeCun, com a LeNet\footfullcite{lecun1989backprop}.
    \item Explodiram em popularidade após 2012, com a AlexNet\footfullcite{krizhevsky2012imagenet}.
    \begin{itemize}
        \item Treinamento em GPUs.
        \item Mais dados.
        \item Mais poder computacional.
    \end{itemize}
\end{itemize}
\end{frame}

\begin{frame}{Problemas com MLP}
\begin{columns}[T]
    \begin{column}{.52\linewidth}
        Imagens têm \textbf{três propriedades} que sugerem o uso de
        arquiteturas específicas:
        \medskip

        \only<1>{%
          \textbf{1.~Alta dimensionalidade}
          \begin{itemize}
            \item Uma imagem $512\times 512$ resulta em $262.144$ pixels.
            \item Com 100 unidades na primeira camada oculta de um MLP,
            teríamos $26.214.400$ pesos.
            \item Em \texttt{float32} (4 bytes por peso), $\approx 100$ MB
            \textbf{só} na primeira camada.
          \end{itemize}}
        \only<2>{%
          \textbf{2.~Alta coesão espacial}
          \begin{itemize}
            \item Pixels próximos têm valores \textbf{altamente
            correlacionados} (veja o \textit{zoom} ao lado).
            \item Uma mesma característica (ex.: uma orelha) pode aparecer
            em \textbf{vários locais} da imagem.
          \end{itemize}}
        \only<3>{%
          \textbf{3.~Semântica estável a transformações geométricas}
          \begin{itemize}
            \item Translações, rotações e deformações \textbf{não mudam}
            o conteúdo: continua sendo um gato.
          \end{itemize}}
    \end{column}
    \begin{column}{.46\linewidth}
        \centering
        \only<1>{%
          \includegraphics[height=5.0cm]{slides/img/02-cnn/mlp_diagram.png}\\[3pt]
          {\scriptsize MLP totalmente conectado: milhões de pesos.}}
        \only<2>{%
          \begin{tikzpicture}[spy using outlines={rectangle, magnification=4,
              size=2.3cm, connect spies}]
            \node[inner sep=0] (dog)
              {\includegraphics[width=4.3cm]{slides/img/02-cnn/mlp_dog.jpeg}};
            \spy[BrickRed, very thick] on (0,-0.55) in node at (3.05,-0.2);
          \end{tikzpicture}\\[3pt]
          {\scriptsize Vizinhança quase uniforme $\to$ valores semelhantes.}}
        \only<3>{%
          \includegraphics[height=2.0cm]{slides/img/02-cnn/mlp_cat.jpeg}\hspace{0.25cm}%
          \rotatebox{15}{\includegraphics[height=2.0cm]{slides/img/02-cnn/mlp_cat.jpeg}}\\[6pt]
          {\scriptsize Original e rotacionado: ainda é ``gato''.}}
    \end{column}
\end{columns}
\end{frame}

\begin{frame}{Camadas Convolucionais}
\begin{block}{Camadas convolucionais trazem diversos benefícios:}
\begin{itemize}
    \item Processam regiões da imagem separadamente, compartilhando pesos.
    \item Usam menos parâmetros do que camadas totalmente conectadas.
    \item Exploram a correlação entre pixels próximos.
    \item Não precisam reaprender padrões em locais diferentes da imagem.
\end{itemize}
\end{block}
\end{frame}

%========================================================================
\section{Invariância e Equivariância}
%========================================================================

\begin{frame}{Invariância}
\begin{itemize}
    \item Uma função $f$ aplicada a uma imagem $\mathbf{X}$ é dita \textbf{invariante} a uma transformação $t$ se:
\end{itemize}
\begin{equation*}
    f\bigl[t[\mathbf{X}]\bigr] = f[\mathbf{X}]
\end{equation*}
\begin{itemize}
    \item Ou seja, independente da transformação aplicada na entrada, o resultado é o mesmo.
    \item Exemplo: classificar a imagem como ``gato'' não deveria depender de o gato estar à esquerda ou à direita do quadro.
\end{itemize}
\end{frame}

\begin{frame}{Equivariância}
\begin{columns}[T]
    \begin{column}{.5\linewidth}
        \begin{itemize}
            \item Uma função $f$ aplicada a uma imagem $\mathbf{X}$ é dita \textbf{equivariante} a uma transformação $t$ se:
        \end{itemize}
        \begin{equation*}
            f\bigl[t[\mathbf{X}]\bigr] = t\bigl[f[\mathbf{X}]\bigr]
        \end{equation*}
        \begin{itemize}
            \item Ou seja, aplicar a transformação na entrada equivale a aplicá-la na saída.
            \item Exemplo: transladar a imagem de entrada deve resultar na mesma segmentação transladada.
        \end{itemize}
    \end{column}
    \begin{column}{.5\linewidth}
        \centering
        \includegraphics[width=\linewidth]{slides/img/02-cnn/equivariance_rooms.png}
    \end{column}
\end{columns}
\end{frame}

%========================================================================
\section{Convolução}
%========================================================================

\begin{frame}{Convolução 2D}
\begin{columns}[T]
    \begin{column}{.5\linewidth}
        \begin{itemize}
            \item A convolução transforma uma imagem de entrada $\mathbf{X}$ em uma imagem de saída $\mathbf{Z}$. Cada posição $(i,j)$ resulta de uma combinação linear de uma janela $3\times 3$ ao redor:
        \end{itemize}
        \begin{equation*}
            z_{i,j} = \sum_{m=1}^{3}\sum_{n=1}^{3} w_{m,n}\, x_{i+m-2,\, j+n-2}
        \end{equation*}
        \begin{itemize}
            \item Os pesos são \textbf{compartilhados} por toda a entrada.
            \item Esse conjunto de pesos é chamado de \textit{kernel} ou \textit{filtro}.
        \end{itemize}
    \end{column}
    \begin{column}{.5\linewidth}
        \centering
        \begin{tikzpicture}[every node/.style={font=\scriptsize}, scale=0.55]
            % Entrada 5x5 com janela 3x3 destacada
            \fill[black!22] (0.5,0.5) rectangle (2.0,2.0);
            \draw[step=0.5, gray!50] (0,0) grid (2.5,2.5);
            \draw[thick, BrickRed] (0.5,0.5) rectangle (2.0,2.0);
            \node[font=\tiny, anchor=south] at (1.25,2.6) {Entrada $5\!\times\!5$};
            % Kernel 3x3 (grid colorido)
            \fill[PastelBlue2!25] (3.5,0.25) rectangle (5.0,1.75);
            \draw[step=0.5, gray!50] (3.5,0.25) grid (5.0,1.75);
            \draw[thick, PastelBlue2!80!black] (3.5,0.25) rectangle (5.0,1.75);
            \node[font=\tiny, anchor=south] at (4.25,1.85) {Kernel};
            % Saida 3x3 com celula destacada
            \fill[black!22] (6.5,1.0) rectangle (7.0,1.5);
            \draw[step=0.5, gray!50] (6.0,0.5) grid (7.5,2.0);
            \draw[thick, BrickRed] (6.5,1.0) rectangle (7.0,1.5);
            \node[font=\tiny, anchor=center] at (6.75,1.25) {$z_{ij}$};
            \node[font=\tiny, anchor=south] at (6.75,2.1) {Saída $3\!\times\!3$};
            % Setas
            \draw[->, BrickRed, thick] (2.55,1.25) -- (3.45,1.0);
            \draw[->, BrickRed, thick] (5.05,1.0) -- (6.45,1.25);
        \end{tikzpicture}
    \end{column}
\end{columns}
\end{frame}

\begin{frame}{Compartilhamento de Pesos}
\begin{columns}[T]
    \begin{column}{.5\linewidth}
        \begin{itemize}
            \item \textbf{O mesmo kernel} $w_{m,n}$ é deslizado por toda a entrada.
            \item O diagrama anterior mostrava apenas \emph{uma} aplicação do kernel; aqui ilustramos a aplicação em duas posições distintas.
            \item Para qualquer posição $(i,j)$:
            \begin{equation*}
                z_{i,j} = \sum_{m,n} w_{m,n}\, x_{i+m-2,\, j+n-2}
            \end{equation*}
            \item Não existem pesos próprios por $x_{i,j}$, os mesmos nove valores reaparecem em cada janela.
        \end{itemize}
    \end{column}
    \begin{column}{.5\linewidth}
        \centering
        \begin{tikzpicture}[every node/.style={font=\scriptsize}, scale=0.55]
            % Entrada 5x5 com duas janelas 3x3 destacadas
            \fill[BrickRed!25] (0.0,2.0) rectangle (1.5,3.5);
            \fill[PastelBlue2!35] (1.0,0.0) rectangle (2.5,1.5);
            \draw[step=0.5, gray!50] (0,0) grid (2.5,3.5);
            \draw[thick, BrickRed] (0.0,2.0) rectangle (1.5,3.5);
            \draw[thick, PastelBlue2!80!black] (1.0,0.0) rectangle (2.5,1.5);
            \node[font=\tiny, anchor=south] at (1.25,3.6) {Entrada $5\!\times\!5$};
            % Kernel (mesmo, grid colorido)
            \fill[PastelBlue2!25] (3.5,1.0) rectangle (5.0,2.5);
            \draw[step=0.5, gray!50] (3.5,1.0) grid (5.0,2.5);
            \draw[thick, PastelBlue2!80!black] (3.5,1.0) rectangle (5.0,2.5);
            \node[font=\tiny, anchor=south] at (4.25,2.6) {Mesmo kernel};
            % Saida 3x3 com duas celulas destacadas
            \fill[BrickRed!25] (6.5,2.0) rectangle (7.0,2.5);
            \fill[PastelBlue2!35] (7.5,0.5) rectangle (8.0,1.0);
            \draw[step=0.5, gray!50] (6.5,0.5) grid (8.0,2.5);
            \draw[thick, BrickRed] (6.5,2.0) rectangle (7.0,2.5);
            \draw[thick, PastelBlue2!80!black] (7.5,0.5) rectangle (8.0,1.0);
            \node[font=\tiny, anchor=south] at (7.25,2.6) {Saída $3\!\times\!3$};
            % Setas
            \draw[->, BrickRed, thick] (1.55,2.75) -- (3.45,2.25);
            \draw[->, BrickRed, thick] (5.05,2.25) -- (6.45,2.25);
            \draw[->, PastelBlue2!80!black, thick] (2.55,0.75) -- (3.45,1.25);
            \draw[->, PastelBlue2!80!black, thick] (5.05,1.25) -- (7.45,0.75);
        \end{tikzpicture}
    \end{column}
\end{columns}
\end{frame}

\begin{frame}{\textit{Padding}}
\begin{columns}[T]
    \begin{column}{.45\linewidth}
        \begin{itemize}
            \item Como lidar com os extremos da entrada?
            \begin{itemize}
                \item \textit{Zero padding}: completar a entrada com zeros.
                \item \textit{Valid padding}: aplicar a convolução apenas onde o \textit{kernel} cabe inteiramente.
                \item Entrada ``circular'' ou periódica (pouco utilizado).
            \end{itemize}
        \end{itemize}
    \end{column}
    \begin{column}{.55\linewidth}
        \centering
        \begin{tikzpicture}[every node/.style={font=\scriptsize}, scale=0.45]
            % Zero padding: entrada 5x5 com borda de zeros (visual 7x7), saida 5x5
            \node[font=\tiny, anchor=south] at (1.75,3.7) {\textit{Zero padding}};
            \fill[black!10] (0,0) rectangle (3.5,3.5);
            \fill[white] (0.5,0.5) rectangle (3.0,3.0);
            \fill[black!22] (0.0,2.0) rectangle (1.5,3.5);
            \draw[step=0.5, gray!50] (0,0) grid (3.5,3.5);
            \draw[thick, BrickRed] (0.0,2.0) rectangle (1.5,3.5);
            % Saida 5x5
            \fill[black!22] (5.0,2.5) rectangle (5.5,3.0);
            \draw[step=0.5, gray!50] (5.0,0.5) grid (7.5,3.0);
            \draw[thick, BrickRed] (5.0,2.5) rectangle (5.5,3.0);
            \node[font=\tiny, anchor=south] at (6.25,3.7) {Saída $5\!\times\!5$};
            \draw[->, BrickRed, thick] (3.55,2.75) -- (4.95,2.75);
            % Valid padding: entrada 5x5 -> saida 3x3
            \node[font=\tiny, anchor=south] at (1.25,-1.1) {\textit{Valid padding}};
            \fill[black!22] (0.0,-3.0) rectangle (1.5,-1.5);
            \draw[step=0.5, gray!50] (0,-4) grid (2.5,-1.5);
            \draw[thick, BrickRed] (0.0,-3.0) rectangle (1.5,-1.5);
            % Saida 3x3
            \fill[black!22] (5.0,-2.5) rectangle (5.5,-2.0);
            \draw[step=0.5, gray!50] (5.0,-3.5) grid (6.5,-2.0);
            \draw[thick, BrickRed] (5.0,-2.5) rectangle (5.5,-2.0);
            \node[font=\tiny, anchor=south] at (5.75,-1.1) {Saída $3\!\times\!3$};
            \draw[->, BrickRed, thick] (2.55,-2.25) -- (4.95,-2.25);
            \node[font=\tiny, align=center, anchor=north] at (3.25,-4.1) {Size$=3$, Stride$=1$, Dilation$=1$};
        \end{tikzpicture}
    \end{column}
\end{columns}
\end{frame}

\begin{frame}{\textit{Stride}}
\begin{columns}[T]
    \begin{column}{.5\linewidth}
        \begin{itemize}
            \item \textit{Stride} é o deslocamento do \textit{kernel} a cada passo.
            \item Pode ser interpretado como \textit{downsampling} quando maior que 1.
            \item Exemplo: Size$=3$, Stride$=2$ produz uma saída com aproximadamente metade da resolução da entrada em cada dimensão.
        \end{itemize}
    \end{column}
    \begin{column}{.5\linewidth}
        \centering
        \begin{tikzpicture}[every node/.style={font=\scriptsize}, scale=0.55]
            % Entrada 5x5 com duas janelas 3x3 destacadas (stride=2)
            \fill[BrickRed!25] (0.0,1.0) rectangle (1.5,2.5);
            \fill[PastelBlue2!35] (1.0,1.0) rectangle (2.5,2.5);
            \draw[step=0.5, gray!50] (0,0) grid (2.5,2.5);
            \draw[thick, BrickRed] (0.0,1.0) rectangle (1.5,2.5);
            \draw[thick, PastelBlue2!80!black, dashed] (1.0,1.0) rectangle (2.5,2.5);
            \node[font=\tiny, anchor=south] at (1.25,2.6) {Entrada $5\!\times\!5$};
            % Saida 2x2
            \fill[BrickRed!25] (4.5,1.5) rectangle (5.0,2.0);
            \fill[PastelBlue2!35] (5.0,1.5) rectangle (5.5,2.0);
            \draw[step=0.5, gray!50] (4.5,1.0) grid (5.5,2.0);
            \draw[thick, BrickRed] (4.5,1.5) rectangle (5.0,2.0);
            \draw[thick, PastelBlue2!80!black, dashed] (5.0,1.5) rectangle (5.5,2.0);
            \node[font=\tiny, anchor=south] at (5.0,2.1) {Saída $2\!\times\!2$};
            % Setas
            \draw[->, BrickRed, thick] (1.55,1.9) -- (4.45,1.85);
            \draw[->, PastelBlue2!80!black, thick, dashed] (2.55,1.6) -- (5.05,1.6);
            \node[font=\tiny, align=center, anchor=north] at (2.75,-0.1) {Size$=3$, Stride$=2$, Dilation$=1$};
        \end{tikzpicture}
    \end{column}
\end{columns}
\end{frame}

\begin{frame}{\textit{Kernel Size}}
\begin{columns}[T]
    \begin{column}{.5\linewidth}
        \begin{itemize}
            \item \textit{Kernel size} é o tamanho do filtro.
            \item Permite integrar informações espaciais mais distantes.
            \begin{itemize}
                \item Mas aumenta o número de pesos a aprender.
            \end{itemize}
            \item Exemplo: Size$=5$, Stride$=1$ olha uma janela $5\times 5$ ao redor de cada posição para produzir cada saída.
        \end{itemize}
    \end{column}
    \begin{column}{.5\linewidth}
        \centering
        \begin{tikzpicture}[every node/.style={font=\scriptsize}, scale=0.5]
            % Entrada 7x7 com janela 5x5 destacada
            \fill[black!22] (0.5,0.5) rectangle (3.0,3.0);
            \draw[step=0.5, gray!50] (0,0) grid (3.5,3.5);
            \draw[thick, BrickRed] (0.5,0.5) rectangle (3.0,3.0);
            \node[font=\tiny, anchor=south] at (1.75,3.6) {Entrada $7\!\times\!7$};
            % Kernel 5x5 (alinhado a grade da entrada: y=0.5..3.0)
            \fill[PastelBlue2!25] (4.5,0.5) rectangle (7.0,3.0);
            \draw[step=0.5, gray!50] (4.5,0.5) grid (7.0,3.0);
            \draw[thick, PastelBlue2!80!black] (4.5,0.5) rectangle (7.0,3.0);
            \node[font=\tiny, anchor=south] at (5.75,3.1) {Kernel $5\!\times\!5$};
            % Saida 3x3 (alinhada: y=1.0..2.5)
            \fill[black!22] (8.0,1.5) rectangle (8.5,2.0);
            \draw[step=0.5, gray!50] (7.5,1.0) grid (9.0,2.5);
            \draw[thick, BrickRed] (8.0,1.5) rectangle (8.5,2.0);
            \node[font=\tiny, anchor=south] at (8.25,2.6) {Saída $3\!\times\!3$};
            % Setas
            \draw[->, BrickRed, thick] (3.55,1.75) -- (4.45,1.75);
            \draw[->, BrickRed, thick] (7.05,1.75) -- (7.95,1.75);
            \node[font=\tiny, align=center, anchor=north] at (4.5,-0.2) {Size$=5$, Stride$=1$, Dilation$=1$};
        \end{tikzpicture}
    \end{column}
\end{columns}
\end{frame}

\begin{frame}{\textit{Dilation}}
\begin{columns}[T]
    \begin{column}{.5\linewidth}
        \begin{itemize}
            \item \textit{Dilation} integra informações mais distantes \textbf{sem} aumento do número de parâmetros.
            \begin{itemize}
                \item Equivalente a aumentar o \textit{kernel} deixando alguns pesos zerados.
            \end{itemize}
            \item Exemplo (convenção dos frameworks como PyTorch): Size$=3$, Stride$=1$, Dilation$=2$ olha posições $(i{\pm}2,\, j{\pm}2)$ em vez de $(i{\pm}1,\, j{\pm}1)$.
            \begin{itemize}
                \item Nessa convenção, Dilation$=1$ corresponde ao kernel denso (sem dilatação) e Dilation$=d$ insere $d-1$ zeros entre os pesos em cada eixo.
            \end{itemize}
        \end{itemize}
    \end{column}
    \begin{column}{.5\linewidth}
        \centering
        \begin{tikzpicture}[every node/.style={font=\scriptsize}, scale=0.55]
            % Entrada 5x5 com 9 posicoes ativas (dilation=2)
            \foreach \i in {0,2,4} {
                \foreach \j in {0,2,4} {
                    \fill[black!22] ({\j*0.5},{\i*0.5}) rectangle ({\j*0.5+0.5},{\i*0.5+0.5});
                }
            }
            \draw[step=0.5, gray!50] (0,0) grid (2.5,2.5);
            \foreach \i in {0,2,4} {
                \foreach \j in {0,2,4} {
                    \draw[thick, BrickRed] ({\j*0.5},{\i*0.5}) rectangle ({\j*0.5+0.5},{\i*0.5+0.5});
                }
            }
            \node[font=\tiny, anchor=south] at (1.25,2.6) {Entrada (área $5\!\times\!5$)};
            % Kernel logico 3x3 (alinhado: y=0.5..2.0)
            \fill[PastelBlue2!25] (3.5,0.5) rectangle (5.0,2.0);
            \draw[step=0.5, gray!50] (3.5,0.5) grid (5.0,2.0);
            \draw[thick, PastelBlue2!80!black] (3.5,0.5) rectangle (5.0,2.0);
            \node[font=\tiny, anchor=south, align=center] at (4.25,2.1) {Kernel $3\!\times\!3$\\dilatado};
            % Saida (alinhada: y=1.0..1.5)
            \fill[black!22] (6.0,1.0) rectangle (6.5,1.5);
            \draw[step=0.5, gray!50] (5.5,0.5) grid (7.0,2.0);
            \draw[thick, BrickRed] (6.0,1.0) rectangle (6.5,1.5);
            \node[font=\tiny, anchor=south] at (6.25,2.1) {Saída};
            % Setas
            \draw[->, BrickRed, thick] (2.55,1.25) -- (3.45,1.25);
            \draw[->, BrickRed, thick] (5.05,1.25) -- (5.95,1.25);
            \node[font=\tiny, align=center, anchor=north] at (4.0,-0.2) {Size$=3$, Stride$=1$, Dilation$=2$};
        \end{tikzpicture}
    \end{column}
\end{columns}
\end{frame}

%========================================================================
\section{Camada Convolucional}
%========================================================================

\begin{frame}{Camada Convolucional}
\begin{itemize}
    \item Uma camada convolucional computa a convolução, soma o \textit{bias} e aplica uma função de ativação:
\end{itemize}
\begin{equation*}
    h_{i,j} = \varphi\!\left[b + \sum_{m=1}^{3}\sum_{n=1}^{3} w_{m,n}\, x_{i+m-2,\, j+n-2}\right]
\end{equation*}
\begin{itemize}
    \item $b$ e os nove pesos $w_{m,n}$ do \textit{kernel} são os parâmetros aprendidos.
    \item A soma cobre toda a janela $K\times K$ ao redor de $(i,j)$.
\end{itemize}
\end{frame}

\begin{frame}{Camada Convolucional vs Totalmente Conectada}
\begin{itemize}
    \item A camada convolucional é um caso especial de camada totalmente conectada.
\end{itemize}
\vspace{0.4em}
\begin{columns}
    \begin{column}{.6\linewidth}
        \centering
        \begin{equation*}
            h_{i,j} = \varphi\!\left[b + \sum_{m=1}^{3}\sum_{n=1}^{3} w_{m,n}\, x_{i+m-2,\, j+n-2}\right]
        \end{equation*}
    \end{column}
    \begin{column}{.4\linewidth}
        \begin{tikzpicture}
            \node[draw, fill=PastelBlue2, text=white, rounded corners, inner sep=4pt] {\small Convolucional};
        \end{tikzpicture}
    \end{column}
\end{columns}
\vspace{0.6em}
\begin{columns}
    \begin{column}{.6\linewidth}
        \centering
        \begin{equation*}
            h_i = \varphi\!\left[b_i + \sum_{j=1}^{D} w_{ij}\, x_{j}\right]
        \end{equation*}
    \end{column}
    \begin{column}{.4\linewidth}
        \begin{tikzpicture}
            \node[draw, fill=PastelBlue2, text=white, rounded corners, inner sep=4pt] {\small Totalmente Conectada};
        \end{tikzpicture}
    \end{column}
\end{columns}
\vspace{0.6em}
\begin{itemize}
    \item Na camada totalmente conectada, $\mathbf{x}$ é a imagem achatada em um vetor de $D$ elementos.
    \item A convolução é uma camada totalmente conectada com pesos compartilhados e a maior parte dos pesos zerados (esparsa e estruturada).
\end{itemize}
\end{frame}

\begin{frame}{Múltiplos Canais e \textit{Feature Maps}}
\begin{itemize}
    \item Tipicamente temos diversos canais de entrada e saída em uma camada convolucional.
    \item As ativações resultantes de um filtro são armazenadas em unidades ocultas tipicamente denominadas \textit{feature maps} (ou canais).
    \item A aplicação de filtros diferentes na mesma entrada gera \textit{feature maps} diferentes.
    \item Em camadas posteriores, cada nova convolução atua sobre todos os \textit{feature maps} simultaneamente.
\end{itemize}
\end{frame}

\begin{frame}{Parâmetros Treináveis}
\begin{itemize}
    \item Considerando $C_i$ canais de entrada e $C_o$ canais de saída, com \textit{kernels} $K\times K$:
\end{itemize}
\begin{equation*}
    \mathbf{\Theta} \in \mathbb{R}^{C_i \times C_o \times K \times K}, \qquad \mathbf{b} \in \mathbb{R}^{C_o}
\end{equation*}
\begin{block}{Exemplo}
Para 3 canais de entrada, 5 canais de saída e \textit{kernels} $3\times 3$ temos $3 \cdot 5 \cdot 3 \cdot 3 = 135$ pesos nos \textit{kernels}, mais 5 \textit{bias}.
\end{block}
\end{frame}

%========================================================================
\section{Campos Receptivos}
%========================================================================

\begin{frame}{Campo Receptivo (\textit{Receptive Field})}
\begin{columns}[T]
    \begin{column}{.5\linewidth}
        \begin{itemize}
            \item Ao realizar uma convolução, cada unidade oculta do \textit{feature map} é gerada levando em conta apenas algumas das entradas.
            \begin{itemize}
                \item Damos o nome de \textbf{Campo Receptivo} à região da entrada considerada por uma unidade oculta.
            \end{itemize}
            \item Para um \textit{kernel} $3\times 3$, cada unidade de $\mathbf{H}_1$ enxerga uma janela $3\times 3$ da entrada.
        \end{itemize}
    \end{column}
    \begin{column}{.5\linewidth}
        \centering
        \begin{tikzpicture}[every node/.style={font=\scriptsize}, scale=0.65]
            % Entrada 5x5 com janela 3x3 destacada
            \fill[black!25] (0.5,1.0) rectangle (2.0,2.5);
            \draw[step=0.5, gray!50] (0,0) grid (2.5,2.5);
            \draw[thick, BrickRed] (0.5,1.0) rectangle (2.0,2.5);
            \node[font=\tiny, anchor=south] at (1.25,2.6) {Entrada, $\mathbf{X}$};
            % H1 (3x3) com uma celula ativa
            \fill[black!25] (4.5,1.5) rectangle (5.0,2.0);
            \draw[step=0.5, gray!50] (4.0,1.0) grid (5.5,2.5);
            \draw[thick, BrickRed] (4.5,1.5) rectangle (5.0,2.0);
            \node[font=\tiny, anchor=south] at (4.75,2.6) {$\mathbf{H}_1$};
            % Cone conectando a janela 3x3 ao ponto ativo
            \fill[gray!25, opacity=0.6] (2.0,2.5) -- (2.0,1.0) -- (4.5,1.75) -- cycle;
            \node[font=\tiny] at (3.3,2.4) {$\Omega\!\in\!\mathbb{R}^{3\times 3}$};
        \end{tikzpicture}
    \end{column}
\end{columns}
\end{frame}

\begin{frame}[shrink=5]{Crescimento do Campo Receptivo}
\begin{columns}[T]
    \begin{column}{.45\linewidth}
        \begin{itemize}
            \item Empilhando convoluções sucessivas, o campo receptivo cresce camada a camada.
            \begin{itemize}
                \item Cada unidade de $\mathbf{H}_2$ enxerga uma janela $3\times 3$ de $\mathbf{H}_1$, que por sua vez vem de uma janela $5\times 5$ na entrada.
                \item Isso permite construir representações cada vez mais ``globais'' da imagem.
            \end{itemize}
            \item Um campo receptivo pequeno pode perder informações relevantes.
            \item Convoluções com \textit{stride} maior que 1 aceleram esse crescimento, ao custo de reduzir a resolução espacial.
        \end{itemize}
    \end{column}
    \begin{column}{.55\linewidth}
        \centering
        \begin{tikzpicture}[every node/.style={font=\scriptsize}, scale=0.55]
            % Entrada 7x7 com janela 5x5 ativa
            \fill[black!25] (0.5,0.5) rectangle (3.0,3.0);
            \draw[step=0.5, gray!50] (0,0) grid (3.5,3.5);
            \draw[thick, BrickRed] (0.5,0.5) rectangle (3.0,3.0);
            \node[font=\tiny, anchor=south] at (1.75,3.6) {$\mathbf{X}$};
            % H1 5x5 com janela 3x3 ativa
            \fill[black!25] (5.0,1.0) rectangle (6.5,2.5);
            \draw[step=0.5, gray!50] (4.5,0.5) grid (7.0,3.0);
            \draw[thick, BrickRed] (5.0,1.0) rectangle (6.5,2.5);
            \node[font=\tiny, anchor=south] at (5.75,3.6) {$\mathbf{H}_1$};
            % H2 3x3 com 1 celula ativa
            \fill[black!25] (8.5,1.5) rectangle (9.0,2.0);
            \draw[step=0.5, gray!50] (8.0,1.0) grid (9.5,2.5);
            \draw[thick, BrickRed] (8.5,1.5) rectangle (9.0,2.0);
            \node[font=\tiny, anchor=south] at (8.75,3.6) {$\mathbf{H}_2$};
            % Cones
            \fill[gray!25, opacity=0.6] (3.0,3.0) -- (3.0,0.5) -- (5.0,1.0) -- (5.0,2.5) -- cycle;
            \fill[gray!25, opacity=0.6] (6.5,2.5) -- (6.5,1.0) -- (8.5,1.75) -- cycle;
            \node[font=\tiny] at (4.0,3.2) {$\Omega\!:\!3\!\times\!3$};
            \node[font=\tiny] at (7.5,2.8) {$\Omega\!:\!3\!\times\!3$};
        \end{tikzpicture}
    \end{column}
\end{columns}
\end{frame}

\begin{frame}{Fórmula do Campo Receptivo}
\begin{itemize}
    \item Para uma sequência de camadas com tamanho de kernel $K_l$ e \textit{stride} $S_l$, o campo receptivo $RF_l$ na camada $l$ obedece à recursão:
\end{itemize}
\begin{equation*}
    RF_l = RF_{l-1} + (K_l - 1) \prod_{i<l} S_i, \qquad RF_0 = 1.
\end{equation*}
\begin{itemize}
    \item Cada nova camada acrescenta $K_l - 1$ posições do espaço original, multiplicadas pelo produto acumulado dos \textit{strides} anteriores.
    \item \textbf{Exemplo}, três convoluções $3\times 3$ com $S=1$:
    \begin{itemize}
        \item $RF_1 = 1 + (3-1)\cdot 1 = 3$.
        \item $RF_2 = 3 + (3-1)\cdot 1 = 5$.
        \item $RF_3 = 5 + (3-1)\cdot 1 = 7$.
    \end{itemize}
    \item Se a segunda camada usa $S_2 = 2$, então $RF_3 = 3 + (3-1)\cdot 1 + (3-1)\cdot 2 = 9$, \textit{stride} acelera o crescimento do campo receptivo.
\end{itemize}
\end{frame}

\begin{frame}{Duas $3\times 3$ ou uma $5\times 5$?}
\begin{columns}[T]
    \begin{column}{.6\linewidth}
        \begin{itemize}
            \item Mesmo campo receptivo, diferentes propriedades.
            \item Resposta: \textbf{duas convoluções $3\times 3$}\footfullcite{simonyan2014vgg}.
            \begin{itemize}
                \item Menos parâmetros: $2 \cdot (3 \cdot 3) = 18$ pesos contra $5 \cdot 5 = 25$.
                \item Mais não linearidade: duas ativações em vez de uma.
            \end{itemize}
        \end{itemize}
    \end{column}
    \begin{column}{.4\linewidth}
        \centering
        \begin{tikzpicture}[scale=0.55, every node/.style={font=\tiny}]
            % Two successive 3x3 convolutions (blue/layered prism)
            % Bottom layer (input, 5x5)
            \def\bx{0}\def\by{0}\def\bw{2.5}\def\bd{1.0}
            \fill[blue!15] (\bx,\by) -- (\bx+\bw,\by) -- (\bx+\bw+\bd,\by+\bd) -- (\bx+\bd,\by+\bd) -- cycle;
            \draw[blue!60] (\bx,\by) -- (\bx+\bw,\by) -- (\bx+\bw+\bd,\by+\bd) -- (\bx+\bd,\by+\bd) -- cycle;
            % Grid lines on bottom layer
            \foreach \gx in {0.6,1.2,1.8,2.4}{ \draw[blue!40, very thin] (\bx+\gx,\by) -- (\bx+\gx+\bd,\by+\bd); }
            \foreach \gy in {0.5,1.0}{ \draw[blue!40, very thin] (\bx,\by+\gy*\bd/1) -- (\bx+\bw,\by+\gy*\bd/1); }
            % Middle layer (3x3 from first conv)
            \def\mx{0.5}\def\my{1.4}
            \fill[blue!25] (\mx,\my) -- (\mx+1.5,\my) -- (\mx+1.5+\bd,\my+\bd) -- (\mx+\bd,\my+\bd) -- cycle;
            \draw[BrickRed!70] (\mx,\my) -- (\mx+1.5,\my) -- (\mx+1.5+\bd,\my+\bd) -- (\mx+\bd,\my+\bd) -- cycle;
            % Top layer (1x1 from second conv)
            \def\tx{1.0}\def\ty{2.8}
            \fill[blue!40] (\tx,\ty) -- (\tx+0.5,\ty) -- (\tx+0.5+\bd,\ty+\bd) -- (\tx+\bd,\ty+\bd) -- cycle;
            \draw[BrickRed!90] (\tx,\ty) -- (\tx+0.5,\ty) -- (\tx+0.5+\bd,\ty+\bd) -- (\tx+\bd,\ty+\bd) -- cycle;
            % Label
            \node[align=center] at (1.8,-0.5) {duas $3\!\times\!3$\\sucessivas};
        \end{tikzpicture}
        \hspace{10pt}
        \begin{tikzpicture}[scale=0.55, every node/.style={font=\tiny}]
            % Single 5x5 convolution (orange prism)
            \def\bx{0}\def\by{0}\def\bw{2.5}\def\bd{1.0}
            \fill[orange!15] (\bx,\by) -- (\bx+\bw,\by) -- (\bx+\bw+\bd,\by+\bd) -- (\bx+\bd,\by+\bd) -- cycle;
            \draw[orange!70] (\bx,\by) -- (\bx+\bw,\by) -- (\bx+\bw+\bd,\by+\bd) -- (\bx+\bd,\by+\bd) -- cycle;
            % Grid lines
            \foreach \gx in {0.5,1.0,1.5,2.0}{ \draw[orange!50, very thin] (\bx+\gx,\by) -- (\bx+\gx+\bd,\by+\bd); }
            \foreach \gy in {0.33,0.66}{ \draw[orange!50, very thin] (\bx,\by) -- (\bx+\bw,\by); }
            % Top layer (1x1)
            \def\tx{1.0}\def\ty{1.4}
            \fill[orange!30] (\tx,\ty) -- (\tx+0.5,\ty) -- (\tx+0.5+\bd,\ty+\bd) -- (\tx+\bd,\ty+\bd) -- cycle;
            \draw[orange!90] (\tx,\ty) -- (\tx+0.5,\ty) -- (\tx+0.5+\bd,\ty+\bd) -- (\tx+\bd,\ty+\bd) -- cycle;
            % Label
            \node[align=center] at (1.8,-0.5) {uma $5\!\times\!5$};
        \end{tikzpicture}
    \end{column}
\end{columns}
\end{frame}

%========================================================================
\section{Treinamento}
%========================================================================

\begin{frame}{Treinamento}
\begin{itemize}
    \item \textit{Backprop} funciona naturalmente para convoluções.
    \begin{itemize}
        \item As derivadas parciais com relação a cada peso compartilhado precisam ser \textbf{acumuladas} (somadas) sobre todas as posições onde esse peso atua.
    \end{itemize}
\end{itemize}
\vspace{0.4em}
\begin{columns}
    \begin{column}{.55\linewidth}
        \footnotesize
        Para um \textit{kernel} $2\times 2$ com pesos $W_{m,n}$ aplicado a uma entrada $X_{i,j}$, a contribuição do gradiente é:
        \begin{align*}
            \partial W_{1,1} &= \sum_{i,j} X_{i,j}\,\partial h_{i,j}\\
            \partial W_{1,2} &= \sum_{i,j} X_{i,j+1}\,\partial h_{i,j}\\
            \partial W_{2,1} &= \sum_{i,j} X_{i+1,j}\,\partial h_{i,j}\\
            \partial W_{2,2} &= \sum_{i,j} X_{i+1,j+1}\,\partial h_{i,j}
        \end{align*}
    \end{column}
    \begin{column}{.45\linewidth}
        \centering
        \begin{tikzpicture}[every node/.style={font=\small}]
            \node[circle, draw, minimum size=24pt] (conv) at (0,0) {CONV};
            \draw[->, thick] (-2.0,0.7) node[left]{$\mathbf{X}$} -- (conv);
            \draw[->, thick] (-2.0,0) node[left]{$\mathbf{\Theta}$} -- (conv);
            \draw[->, thick] (-2.0,-0.7) node[left]{$\mathbf{b}$} -- (conv);
            \draw[->, thick] (conv) -- (1.8,0) node[right]{$\mathbf{Y}$};
            \node[BrickRed, font=\scriptsize] at (-1.6,0.95) {$\partial \mathbf{X}$};
            \node[BrickRed, font=\scriptsize] at (-1.6,0.25) {$\partial \mathbf{\Theta}$};
            \node[BrickRed, font=\scriptsize] at (-1.6,-0.45) {$\partial \mathbf{b}$};
            \node[BrickRed, font=\scriptsize] at (1.4,0.25) {$\partial \mathbf{Y}$};
        \end{tikzpicture}
    \end{column}
\end{columns}
\end{frame}

\begin{frame}{O Resultado do Treinamento}
\begin{itemize}
    \item Os \textit{feature maps} das primeiras camadas mostram arestas, cantos e formas primitivas.
    \item Em camadas mais internas vemos formas mais complexas: olhos, narizes, partes de objetos.
    \begin{itemize}
        \item Repare que aumenta também o campo receptivo.
    \end{itemize}
    \item Esse comportamento hierárquico não é programado: emerge naturalmente do treinamento.
\end{itemize}
\end{frame}

%========================================================================
\section{Downsampling e Upsampling}
%========================================================================

\begin{frame}{\textit{Downsampling}}
\begin{columns}[T]
    \begin{column}{.55\linewidth}
        \begin{itemize}
            \item Subamostrar espacialmente a imagem pode ter o efeito de aumentar o campo receptivo de uma unidade oculta.
            \item As principais formas de subamostragem em CNN são:
            \begin{enumerate}
                \item Usar \textit{stride} maior que 1 na convolução.
                \item \textit{Max pooling}: mantém apenas o maior valor de cada janela.
                \item \textit{Average pooling}: usa a média dos valores de cada janela.
            \end{enumerate}
        \end{itemize}
    \end{column}
    \begin{column}{.45\linewidth}
        \centering
        \begin{tikzpicture}[every node/.style={font=\scriptsize}, scale=0.42]
            % --- Max pooling ---
            \fill[PastelBlue2!40] (0,2) rectangle (2,4);
            \fill[orange!30]      (2,2) rectangle (4,4);
            \fill[green!25]       (0,0) rectangle (2,2);
            \fill[BrickRed!25]    (2,0) rectangle (4,2);
            \draw[gray!60] (0,0) rectangle (4,4);
            \draw[gray!60] (1,0) -- (1,4);
            \draw[gray!60] (2,0) -- (2,4);
            \draw[gray!60] (3,0) -- (3,4);
            \draw[gray!60] (0,1) -- (4,1);
            \draw[gray!60] (0,2) -- (4,2);
            \draw[gray!60] (0,3) -- (4,3);
            \draw[thick, PastelBlue2!70!black] (0,2) rectangle (2,4);
            \draw[thick, orange!70!black]      (2,2) rectangle (4,4);
            \draw[thick, green!50!black]       (0,0) rectangle (2,2);
            \draw[thick, BrickRed!80]          (2,0) rectangle (4,2);
            \node at (0.5,3.5) {$1$};   \node at (1.5,3.5) {$3$};
            \node at (0.5,2.5) {$2$};   \node at (1.5,2.5) {$\mathbf{8}$};
            \node at (2.5,3.5) {$\mathbf{9}$};   \node at (3.5,3.5) {$4$};
            \node at (2.5,2.5) {$1$};   \node at (3.5,2.5) {$5$};
            \node at (0.5,1.5) {$4$};   \node at (1.5,1.5) {$2$};
            \node at (0.5,0.5) {$\mathbf{7}$};   \node at (1.5,0.5) {$1$};
            \node at (2.5,1.5) {$3$};   \node at (3.5,1.5) {$\mathbf{6}$};
            \node at (2.5,0.5) {$2$};   \node at (3.5,0.5) {$5$};
            \node[font=\tiny, anchor=south] at (2,4.05) {entrada};
            \draw[->, gray!70, thick] (4.2,2) -- (5.4,2);
            \node[font=\tiny, anchor=south] at (4.8,2.1) {max};
            \fill[PastelBlue2!40] (5.7,3) rectangle (6.7,4);
            \fill[orange!30]      (6.7,3) rectangle (7.7,4);
            \fill[green!25]       (5.7,2) rectangle (6.7,3);
            \fill[BrickRed!25]    (6.7,2) rectangle (7.7,3);
            \draw[gray!60] (5.7,2) rectangle (7.7,4);
            \draw[gray!60] (6.7,2) -- (6.7,4);
            \draw[gray!60] (5.7,3) -- (7.7,3);
            \node at (6.2,3.5) {$8$};
            \node at (7.2,3.5) {$9$};
            \node at (6.2,2.5) {$7$};
            \node at (7.2,2.5) {$6$};
            \node[font=\tiny, anchor=south] at (6.7,4.05) {saída};

            % --- Average pooling, deslocado para baixo ---
            \begin{scope}[yshift=-5cm]
            \fill[PastelBlue2!40] (0,2) rectangle (2,4);
            \fill[orange!30]      (2,2) rectangle (4,4);
            \fill[green!25]       (0,0) rectangle (2,2);
            \fill[BrickRed!25]    (2,0) rectangle (4,2);
            \draw[gray!60] (0,0) rectangle (4,4);
            \draw[gray!60] (1,0) -- (1,4);
            \draw[gray!60] (2,0) -- (2,4);
            \draw[gray!60] (3,0) -- (3,4);
            \draw[gray!60] (0,1) -- (4,1);
            \draw[gray!60] (0,2) -- (4,2);
            \draw[gray!60] (0,3) -- (4,3);
            \draw[thick, PastelBlue2!70!black] (0,2) rectangle (2,4);
            \draw[thick, orange!70!black]      (2,2) rectangle (4,4);
            \draw[thick, green!50!black]       (0,0) rectangle (2,2);
            \draw[thick, BrickRed!80]          (2,0) rectangle (4,2);
            \node at (0.5,3.5) {$2$};   \node at (1.5,3.5) {$4$};
            \node at (0.5,2.5) {$0$};   \node at (1.5,2.5) {$2$};
            \node at (2.5,3.5) {$6$};   \node at (3.5,3.5) {$8$};
            \node at (2.5,2.5) {$4$};   \node at (3.5,2.5) {$6$};
            \node at (0.5,1.5) {$1$};   \node at (1.5,1.5) {$3$};
            \node at (0.5,0.5) {$3$};   \node at (1.5,0.5) {$5$};
            \node at (2.5,1.5) {$5$};   \node at (3.5,1.5) {$7$};
            \node at (2.5,0.5) {$7$};   \node at (3.5,0.5) {$9$};
            \node[font=\tiny, anchor=south] at (2,4.05) {entrada};
            \draw[->, gray!70, thick] (4.2,2) -- (5.4,2);
            \node[font=\tiny, anchor=south] at (4.8,2.1) {avg};
            \fill[PastelBlue2!40] (5.7,3) rectangle (6.7,4);
            \fill[orange!30]      (6.7,3) rectangle (7.7,4);
            \fill[green!25]       (5.7,2) rectangle (6.7,3);
            \fill[BrickRed!25]    (6.7,2) rectangle (7.7,3);
            \draw[gray!60] (5.7,2) rectangle (7.7,4);
            \draw[gray!60] (6.7,2) -- (6.7,4);
            \draw[gray!60] (5.7,3) -- (7.7,3);
            \node at (6.2,3.5) {$2$};
            \node at (7.2,3.5) {$6$};
            \node at (6.2,2.5) {$3$};
            \node at (7.2,2.5) {$7$};
            \node[font=\tiny, anchor=south] at (6.7,4.05) {saída};
            \end{scope}
        \end{tikzpicture}
    \end{column}
\end{columns}
\end{frame}

\begin{frame}{\textit{Upsampling}}
\begin{columns}[T]
    \begin{column}{.55\linewidth}
        \begin{itemize}
            \item Sobreamostrar espacialmente pode ser útil quando a saída da rede também é uma imagem (segmentação, super resolução, geração).
            \item As principais formas de sobreamostragem em CNN são:
            \begin{enumerate}
                \item Duplicar os valores conhecidos.
                \item \textit{Max unpooling}: insere o valor na posição original do \textit{max pooling} e zera as demais.
                \item Interpolação bilinear.
                \item Convolução transposta.
            \end{enumerate}
        \end{itemize}
    \end{column}
    \begin{column}{.45\linewidth}
        \centering
        \begin{tikzpicture}[every node/.style={font=\scriptsize}, scale=0.42]
            % --- Duplicação de valores (nearest neighbor) ---
            \fill[PastelBlue2!40] (0,2) rectangle (1,3);
            \fill[orange!30]      (1,2) rectangle (2,3);
            \fill[green!25]       (0,1) rectangle (1,2);
            \fill[BrickRed!25]    (1,1) rectangle (2,2);
            \draw[gray!60] (0,1) rectangle (2,3);
            \draw[gray!60] (1,1) -- (1,3);
            \draw[gray!60] (0,2) -- (2,2);
            \node at (0.5,2.5) {$8$};
            \node at (1.5,2.5) {$9$};
            \node at (0.5,1.5) {$7$};
            \node at (1.5,1.5) {$6$};
            \node[font=\tiny, anchor=south] at (1,3.05) {entrada};
            \draw[->, gray!70, thick] (2.2,2) -- (5.2,2);
            \node[font=\tiny, anchor=south] at (3.7,2.1) {duplicar};
            \fill[PastelBlue2!40] (5.5,2) rectangle (7.5,4);
            \fill[orange!30]      (7.5,2) rectangle (9.5,4);
            \fill[green!25]       (5.5,0) rectangle (7.5,2);
            \fill[BrickRed!25]    (7.5,0) rectangle (9.5,2);
            \draw[gray!60] (5.5,0) rectangle (9.5,4);
            \draw[gray!60] (6.5,0) -- (6.5,4);
            \draw[gray!60] (7.5,0) -- (7.5,4);
            \draw[gray!60] (8.5,0) -- (8.5,4);
            \draw[gray!60] (5.5,1) -- (9.5,1);
            \draw[gray!60] (5.5,2) -- (9.5,2);
            \draw[gray!60] (5.5,3) -- (9.5,3);
            \node at (6.0,3.5) {$8$}; \node at (7.0,3.5) {$8$};
            \node at (8.0,3.5) {$9$}; \node at (9.0,3.5) {$9$};
            \node at (6.0,2.5) {$8$}; \node at (7.0,2.5) {$8$};
            \node at (8.0,2.5) {$9$}; \node at (9.0,2.5) {$9$};
            \node at (6.0,1.5) {$7$}; \node at (7.0,1.5) {$7$};
            \node at (8.0,1.5) {$6$}; \node at (9.0,1.5) {$6$};
            \node at (6.0,0.5) {$7$}; \node at (7.0,0.5) {$7$};
            \node at (8.0,0.5) {$6$}; \node at (9.0,0.5) {$6$};
            \node[font=\tiny, anchor=south] at (7.5,4.05) {saída};

            % --- Interpolação bilinear, deslocada para baixo ---
            \begin{scope}[yshift=-5cm]
            % Entrada 2x2
            \fill[PastelBlue2!40] (0,2) rectangle (1,3);
            \fill[orange!30]      (1,2) rectangle (2,3);
            \fill[green!25]       (0,1) rectangle (1,2);
            \fill[BrickRed!25]    (1,1) rectangle (2,2);
            \draw[gray!60] (0,1) rectangle (2,3);
            \draw[gray!60] (1,1) -- (1,3);
            \draw[gray!60] (0,2) -- (2,2);
            \node at (0.5,2.5) {$0$};
            \node at (1.5,2.5) {$3$};
            \node at (0.5,1.5) {$6$};
            \node at (1.5,1.5) {$9$};
            \node[font=\tiny, anchor=south] at (1,3.05) {entrada};
            \draw[->, gray!70, thick] (2.2,2) -- (5.2,2);
            \node[font=\tiny, anchor=south] at (3.7,2.1) {bilinear};
            % Saída 4x4 (valores interpolados, cantos mantêm os 4 originais)
            \fill[PastelBlue2!40] (5.5,2) rectangle (7.5,4);
            \fill[orange!30]      (7.5,2) rectangle (9.5,4);
            \fill[green!25]       (5.5,0) rectangle (7.5,2);
            \fill[BrickRed!25]    (7.5,0) rectangle (9.5,2);
            \draw[gray!60] (5.5,0) rectangle (9.5,4);
            \draw[gray!60] (6.5,0) -- (6.5,4);
            \draw[gray!60] (7.5,0) -- (7.5,4);
            \draw[gray!60] (8.5,0) -- (8.5,4);
            \draw[gray!60] (5.5,1) -- (9.5,1);
            \draw[gray!60] (5.5,2) -- (9.5,2);
            \draw[gray!60] (5.5,3) -- (9.5,3);
            \node at (6.0,3.5) {$0$}; \node at (7.0,3.5) {$1$};
            \node at (8.0,3.5) {$2$}; \node at (9.0,3.5) {$3$};
            \node at (6.0,2.5) {$2$}; \node at (7.0,2.5) {$3$};
            \node at (8.0,2.5) {$4$}; \node at (9.0,2.5) {$5$};
            \node at (6.0,1.5) {$4$}; \node at (7.0,1.5) {$5$};
            \node at (8.0,1.5) {$6$}; \node at (9.0,1.5) {$7$};
            \node at (6.0,0.5) {$6$}; \node at (7.0,0.5) {$7$};
            \node at (8.0,0.5) {$8$}; \node at (9.0,0.5) {$9$};
            \node[font=\tiny, anchor=south] at (7.5,4.05) {saída};
            \end{scope}
        \end{tikzpicture}
    \end{column}
\end{columns}
\end{frame}

\begin{frame}{Convolução Transposta}
\begin{itemize}
    \item Operação aprendível para aumentar a resolução espacial.
    \item Cada elemento da entrada multiplica o \textit{kernel} inteiro e o resultado é depositado em uma região da saída, com sobreposições somadas.
    \item Parâmetros: \textit{kernel size}, \textit{padding}, \textit{stride}. \textit{Stride} maior que 1 \textbf{expande} a saída.
\end{itemize}
\end{frame}

\begin{frame}{Convolução Transposta: Exemplo}
\begin{itemize}
    \item Filtro $3\times 3$, \textit{padding} $=1$, \textit{stride} $=2$. Cada célula da entrada projeta o \textit{kernel} inteiro multiplicado por seu valor; sobreposições são somadas.
\end{itemize}
\vspace{0.2em}
\centering
\begin{tikzpicture}[every node/.style={font=\scriptsize}, scale=0.62,
    cory/.style={fill=orange!40},
    cpbb/.style={fill=PastelBlue2!45},
    cgrb/.style={fill=green!35},
    cbrb/.style={fill=BrickRed!35},
    tory/.style={text=orange!70!black},
    tpbb/.style={text=PastelBlue2!60!black},
    tgrb/.style={text=green!45!black},
    tbrb/.style={text=BrickRed},
    every cellnode/.style={font=\footnotesize}
]
    % Entrada 2x2
    \fill[cory] (0,1) rectangle (1,2);
    \fill[cpbb] (1,1) rectangle (2,2);
    \fill[cgrb] (0,0) rectangle (1,1);
    \fill[cbrb] (1,0) rectangle (2,1);
    \draw[step=1.0, gray!60] (0,0) grid (2,2);
    \node at (0.5,1.5) {$2$}; \node at (1.5,1.5) {$1$};
    \node at (0.5,0.5) {$3$}; \node at (1.5,0.5) {$2$};
    \node[font=\tiny, anchor=south] at (1,2.05) {entrada $\mathbf{X}$};
    \only<1>{\draw[ultra thick, purple] (0,1) rectangle (1,2);}
    \only<2>{\draw[ultra thick, purple] (1,1) rectangle (2,2);}
    \only<3>{\draw[ultra thick, purple] (0,0) rectangle (1,1);}
    \only<4>{\draw[ultra thick, purple] (1,0) rectangle (2,1);}

    % Kernel 3x3
    \fill[PastelBlue2!15] (3.0,-0.5) rectangle (6.0,2.5);
    \draw[gray!60] (3.0,-0.5) rectangle (6.0,2.5);
    \draw[gray!60] (4.0,-0.5) -- (4.0,2.5);
    \draw[gray!60] (5.0,-0.5) -- (5.0,2.5);
    \draw[gray!60] (3.0,0.5) -- (6.0,0.5);
    \draw[gray!60] (3.0,1.5) -- (6.0,1.5);
    \node at (3.5,2.0) {$1$}; \node at (4.5,2.0) {$2$}; \node at (5.5,2.0) {$1$};
    \node at (3.5,1.0) {$2$}; \node at (4.5,1.0) {$0$}; \node at (5.5,1.0) {$1$};
    \node at (3.5,0.0) {$0$}; \node at (4.5,0.0) {$2$}; \node at (5.5,0.0) {$1$};
    \node[font=\tiny, anchor=south] at (4.5,2.55) {kernel $\mathbf{\Theta}$};

    % Resultado 4x4 (com padding-1 visualmente sombreado fora)
    \draw[step=1.0, gray!60] (7.0,-1.0) grid (11.0,3.0);
    \node[font=\tiny, anchor=south] at (9.0,3.05) {resultado $4\!\times\!4$};

    % --- Conteudo das celulas (um node por celula, conteudo cumulativo) ---
    % Cada celula compoe contribuicoes de input cells em ordem: yellow, blue, green, red
    % Row 1 (top, y=2.5)
    \node[font=\tiny] at (7.5,2.5) {\only<1->{\textcolor{orange!70!black}{$0$}}};
    \node[font=\tiny] at (8.5,2.5) {\only<1->{\textcolor{orange!70!black}{$2$}}\only<2->{\textcolor{PastelBlue2!60!black}{$\!+\!2$}}};
    \node[font=\tiny] at (9.5,2.5) {\only<2->{\textcolor{PastelBlue2!60!black}{$0$}}};
    \node[font=\tiny] at (10.5,2.5) {\only<2->{\textcolor{PastelBlue2!60!black}{$1$}}};
    % Row 2 (y=1.5)
    \node[font=\tiny] at (7.5,1.5) {\only<1->{\textcolor{orange!70!black}{$4$}}\only<3->{\textcolor{green!45!black}{$\!+\!6$}}};
    \node[font=\tiny] at (8.5,1.5) {\only<1->{\textcolor{orange!70!black}{$2$}}\only<2->{\textcolor{PastelBlue2!60!black}{$\!+\!0$}}\only<3->{\textcolor{green!45!black}{$\!+\!3$}}\only<4->{\textcolor{BrickRed}{$\!+\!2$}}};
    \node[font=\tiny] at (9.5,1.5) {\only<2->{\textcolor{PastelBlue2!60!black}{$2$}}\only<4->{\textcolor{BrickRed}{$\!+\!4$}}};
    \node[font=\tiny] at (10.5,1.5) {\only<2->{\textcolor{PastelBlue2!60!black}{$1$}}\only<4->{\textcolor{BrickRed}{$\!+\!2$}}};
    % Row 3 (y=0.5)
    \node[font=\tiny] at (7.5,0.5) {\only<3->{\textcolor{green!45!black}{$0$}}};
    \node[font=\tiny] at (8.5,0.5) {\only<3->{\textcolor{green!45!black}{$3$}}\only<4->{\textcolor{BrickRed}{$\!+\!4$}}};
    \node[font=\tiny] at (9.5,0.5) {\only<4->{\textcolor{BrickRed}{$0$}}};
    \node[font=\tiny] at (10.5,0.5) {\only<4->{\textcolor{BrickRed}{$2$}}};
    % Row 4 (y=-0.5)
    \node[font=\tiny] at (7.5,-0.5) {\only<3->{\textcolor{green!45!black}{$6$}}};
    \node[font=\tiny] at (8.5,-0.5) {\only<3->{\textcolor{green!45!black}{$3$}}\only<4->{\textcolor{BrickRed}{$\!+\!0$}}};
    \node[font=\tiny] at (9.5,-0.5) {\only<4->{\textcolor{BrickRed}{$4$}}};
    \node[font=\tiny] at (10.5,-0.5) {\only<4->{\textcolor{BrickRed}{$2$}}};
\end{tikzpicture}
\end{frame}

\begin{frame}{Convolução $1\times 1$}
\begin{columns}[T]
    \begin{column}{.5\linewidth}
        \begin{itemize}
            \item Algumas vezes queremos mudar o número de canais sem agrupamento espacial.
            \begin{itemize}
                \item Nesses casos podemos aplicar a convolução $1\times 1$.
            \end{itemize}
            \item Cada posição espacial é independente: cada saída $h'_{i,j,c}$ é uma combinação linear dos canais de entrada $h_{i,j,1\dots C_i}$.
            \item Útil para reduzir ou expandir a profundidade do tensor sem alterar largura e altura.
        \end{itemize}
    \end{column}
    \begin{column}{.5\linewidth}
        \centering
        \begin{tikzpicture}[every node/.style={font=\scriptsize}, scale=0.5]
            % 3 canais de entrada empilhados (offset diagonal)
            \foreach \k/\clr/\dx/\dy in {3/PastelBlue2!15/0.5/0.5, 2/PastelBlue2!25/0.25/0.25, 1/PastelBlue2!40/0/0} {
                \fill[\clr] (\dx,\dy) rectangle (\dx+2.0,\dy+2.0);
                \draw[step=0.5, gray!50] (\dx,\dy) rectangle (\dx+2.0,\dy+2.0);
                \draw[gray!50] (\dx+0.5,\dy) -- (\dx+0.5,\dy+2.0);
                \draw[gray!50] (\dx+1.0,\dy) -- (\dx+1.0,\dy+2.0);
                \draw[gray!50] (\dx+1.5,\dy) -- (\dx+1.5,\dy+2.0);
                \draw[gray!50] (\dx,\dy+0.5) -- (\dx+2.0,\dy+0.5);
                \draw[gray!50] (\dx,\dy+1.0) -- (\dx+2.0,\dy+1.0);
                \draw[gray!50] (\dx,\dy+1.5) -- (\dx+2.0,\dy+1.5);
                \draw[thick, PastelBlue2!80!black] (\dx,\dy) rectangle (\dx+2.0,\dy+2.0);
            }
            \node[font=\tiny, anchor=south] at (1.25,3.0) {entrada (3 canais)};

            % Kernel 1x1 x 3 (coluna estreita de 3 celulas, cores empilhadas, centralizada no eixo do stack)
            \fill[BrickRed!45] (4.0,1.75) rectangle (4.5,2.25);
            \fill[BrickRed!30] (4.0,1.25) rectangle (4.5,1.75);
            \fill[BrickRed!15] (4.0,0.75) rectangle (4.5,1.25);
            \draw[gray!50] (4.0,0.75) rectangle (4.5,2.25);
            \draw[gray!50] (4.0,1.25) -- (4.5,1.25);
            \draw[gray!50] (4.0,1.75) -- (4.5,1.75);
            \draw[thick, BrickRed] (4.0,0.75) rectangle (4.5,2.25);
            \node[font=\tiny, anchor=north, align=center] at (4.25,0.70) {kernel\\$1\!\times\!1\!\times\!3$};

            % Canal de saida (1 canal)
            \fill[orange!25] (5.7,0.5) rectangle (7.7,2.5);
            \draw[step=0.5, gray!50] (5.7,0.5) grid (7.7,2.5);
            \draw[thick, orange!70!black] (5.7,0.5) rectangle (7.7,2.5);
            \node[font=\tiny, anchor=south] at (6.7,2.55) {saída (1 canal)};

            % Setas
            \draw[->, gray!70, thick] (3.05,1.5) -- (3.95,1.5);
            \draw[->, gray!70, thick] (4.55,1.5) -- (5.65,1.5);
        \end{tikzpicture}
    \end{column}
\end{columns}
\end{frame}

%========================================================================
\section{Implementação de Convolução}
%========================================================================

\begin{frame}[fragile]{Implementação: Versão Ingênua}
\begin{itemize}
    \item Itera explicitamente por linhas, colunas, canais e \textit{kernel}:
\end{itemize}
\begin{lstlisting}[style=vscode-light, basicstyle=\ttfamily\scriptsize, language=Python]
# Versao com um unico canal de saida (C_o = 1).
# Para multiplos canais de saida bastaria envolver tudo
# em mais um loop sobre co e indexar feat_map_out[m, n, co].
for m in range(0, M):           # linhas
    for n in range(0, N):       # colunas
        for c in range(0, C):   # canais de entrada
            for i in range(0, K):     # tamanho kernel
                for j in range(0, K): # tamanho kernel
                    feat_map_out[m, n] += kernel[i, j, c] * feat_map_in[m+i, n+j, c]
        feat_map_out[m, n] += bias
\end{lstlisting}
\begin{itemize}
    \item Simples de entender, mas extremamente lenta e sem aproveitar paralelismo de hardware.
\end{itemize}
\end{frame}

\begin{frame}{Implementação: via Transformada de Fourier}
\begin{itemize}
    \item Versão complexa e raramente utilizada na prática.
    \item Convolução no domínio espacial é equivalente a multiplicação no domínio da frequência:
\end{itemize}
\begin{equation*}
    g(x) = f(x) * h(x) \qquad \Longleftrightarrow \qquad G(u) = F(u)\,H(u)
\end{equation*}
\vspace{-0.4em}
\begin{center}
\begin{tikzpicture}[every node/.style={font=\scriptsize},
    box/.style={draw, rounded corners=2pt, minimum width=1.6cm, minimum height=0.55cm, inner sep=1pt},
    sig/.style={draw, rounded corners=2pt, minimum width=0.9cm, minimum height=0.55cm, inner sep=1pt},
    op/.style={draw, circle, minimum size=0.45cm, inner sep=0pt},
    >=stealth
]
    \node[sig, fill=pink!50] (f) at (0,0.6) {$f(x)$};
    \node[sig, fill=yellow!40] (h) at (0,-0.6) {$h(x)$};
    \node[box, fill=gray!20] (ff) at (1.8,0.6) {FFT};
    \node[box, fill=gray!20] (fh) at (1.8,-0.6) {FFT};
    \node[sig, fill=pink!80] (F) at (3.6,0.6) {$F(u)$};
    \node[sig, fill=yellow!80] (H) at (3.6,-0.6) {$H(u)$};
    \node[op, fill=PastelBlue2!30] (mul) at (5.0,0) {$\times$};
    \node[box, fill=gray!20] (ifft) at (6.8,0) {IFFT};
    \node[sig, fill=green!40, minimum width=1.2cm] (out) at (8.6,0) {$f * h$};
    \draw[->] (f) -- (ff);
    \draw[->] (h) -- (fh);
    \draw[->] (ff) -- (F);
    \draw[->] (fh) -- (H);
    \draw[->] (F.east) -- (mul);
    \draw[->] (H.east) -- (mul);
    \draw[->] (mul) -- (ifft);
    \draw[->] (ifft) -- (out);
\end{tikzpicture}
\end{center}
\begin{itemize}
    \item Útil para convoluções com \textit{kernels} grandes (acima de $11\times 11$)\footfullcite{vasilache2014fbfft}.
\end{itemize}
\end{frame}

\begin{frame}{Implementação: \texttt{im2col} e Matriz de Toeplitz}
\begin{itemize}
    \item Versão utilizada na prática em GPUs e bibliotecas otimizadas.
    \item Existe uma forma inteligente de organizar a matriz de pesos para obter o resultado da convolução com uma única multiplicação de matrizes.
    \item Uma \textbf{matriz de Toeplitz} é uma matriz cujas diagonais principal e subdiagonais são constantes:
\end{itemize}
\begin{equation*}
\begin{bmatrix}
5 & 27 & 2 & -56 & 0 & 1 \\
-3 & 5 & 27 & 2 & -56 & 0 \\
8 & -3 & 5 & 27 & 2 & -56 \\
1 & 8 & -3 & 5 & 27 & 2 \\
-9 & 1 & 8 & -3 & 5 & 27 \\
3 & -9 & 1 & 8 & -3 & 5
\end{bmatrix}
\end{equation*}
\end{frame}

\begin{frame}{\texttt{im2col}: Passo 1 -- Entrada e filtro}
\begin{itemize}
    \item Definimos a entrada $\mathbf{X}$ ($3\times 3$) e o filtro $\mathbf{\Theta}$ ($2\times 2$):
\end{itemize}
\vspace{0.4em}
\centering
\begin{tikzpicture}[every node/.style={font=\footnotesize}, scale=0.8]
    % X 3x3 azul
    \node at (-0.7,1.5) {$\mathbf{X} = $};
    \fill[PastelBlue2!35] (0,0) rectangle (3,3);
    \draw[step=1, gray!60] (0,0) grid (3,3);
    \node at (0.5,2.5) {$x_{1,1}$}; \node at (1.5,2.5) {$x_{1,2}$}; \node at (2.5,2.5) {$x_{1,3}$};
    \node at (0.5,1.5) {$x_{2,1}$}; \node at (1.5,1.5) {$x_{2,2}$}; \node at (2.5,1.5) {$x_{2,3}$};
    \node at (0.5,0.5) {$x_{3,1}$}; \node at (1.5,0.5) {$x_{3,2}$}; \node at (2.5,0.5) {$x_{3,3}$};
    % w 2x2 laranja
    \node at (5.3,1.5) {$\mathbf{\Theta} = $};
    \fill[orange!40] (6.0,0.5) rectangle (8.0,2.5);
    \draw[gray!60] (6.0,0.5) rectangle (8.0,2.5);
    \draw[gray!60] (7.0,0.5) -- (7.0,2.5);
    \draw[gray!60] (6.0,1.5) -- (8.0,1.5);
    \node at (6.5,2.0) {$w_{1,1}$}; \node at (7.5,2.0) {$w_{1,2}$};
    \node at (6.5,1.0) {$w_{2,1}$}; \node at (7.5,1.0) {$w_{2,2}$};
\end{tikzpicture}
\end{frame}

\begin{frame}{\texttt{im2col}: Passo 2 -- Tamanho da saída}
\begin{itemize}
    \item Modo \textit{full}: $(X_h + w_h - 1)\times(X_w + w_w - 1)$. Para $\mathbf{X}$ $3\times 3$ e $\mathbf{\Theta}$ $2\times 2$:
\end{itemize}
\begin{equation*}
    (3 + 2 - 1) \times (3 + 2 - 1) = 4 \times 4
\end{equation*}
\begin{itemize}
    \item Convenções de tamanho:
    \begin{itemize}
        \item \textbf{full} (\textit{signal processing}): saída maior que a entrada; considera todas as janelas em que entrada e kernel se sobrepõem em pelo menos um pixel.
        \item \textbf{valid} (típico em CNNs): saída $(X_h - w_h + 1) \times (X_w - w_w + 1)$; considera apenas janelas inteiramente dentro da entrada.
        \item \textbf{same}: com \textit{padding} adequado, mantém o tamanho original da entrada.
    \end{itemize}
\end{itemize}
\end{frame}

\begin{frame}{\texttt{im2col}: Passo 3 -- Zero padding do filtro}
\begin{itemize}
    \item Aplicamos \textit{zero padding} no filtro até atingir o tamanho da saída $4\times 4$:
\end{itemize}
\vspace{0.4em}
\centering
\begin{tikzpicture}[every node/.style={font=\footnotesize}, scale=0.8]
    \node at (-0.9,2.0) {$\mathbf{\Theta} = $};
    % Padded 4x4: zeros no topo e a direita, pesos no canto inferior-esquerdo
    \fill[gray!15] (0,0) rectangle (4,4);
    \fill[orange!40] (0,0) rectangle (2,2);
    \draw[step=1, gray!60] (0,0) grid (4,4);
    \node at (0.5,3.5) {$0$}; \node at (1.5,3.5) {$0$}; \node at (2.5,3.5) {$0$}; \node at (3.5,3.5) {$0$};
    \node at (0.5,2.5) {$0$}; \node at (1.5,2.5) {$0$}; \node at (2.5,2.5) {$0$}; \node at (3.5,2.5) {$0$};
    \node at (0.5,1.5) {$w_{1,1}$}; \node at (1.5,1.5) {$w_{1,2}$};
    \node at (2.5,1.5) {$0$};       \node at (3.5,1.5) {$0$};
    \node at (0.5,0.5) {$w_{2,1}$}; \node at (1.5,0.5) {$w_{2,2}$};
    \node at (2.5,0.5) {$0$};       \node at (3.5,0.5) {$0$};
\end{tikzpicture}
\end{frame}

\begin{frame}{\texttt{im2col}: Passo 4 -- Matrizes de Toeplitz por linha}
\begin{itemize}
    \item Para cada linha do filtro paddado, construímos uma matriz de Toeplitz $\mathbf{F}_k$ (tamanho saída$_h \times$ entrada$_w = 4\times 3$). Começamos pela última linha.
\end{itemize}
\vspace{0.3em}
\centering
\begin{tikzpicture}[every node/.style={font=\tiny}, scale=0.55]
    % F0 azul (da linha 4 do padded: w_{2,1}, w_{2,2}, 0, 0)
    \node[anchor=south, font=\scriptsize] at (1.5,4.2) {$\mathbf{F}_0$};
    \fill[PastelBlue2!35] (0,0) rectangle (3,4);
    \draw[step=1, gray!60] (0,0) grid (3,4);
    \node at (0.5,3.5) {$w_{2,1}$}; \node at (1.5,3.5) {$0$}; \node at (2.5,3.5) {$0$};
    \node at (0.5,2.5) {$w_{2,2}$}; \node at (1.5,2.5) {$w_{2,1}$}; \node at (2.5,2.5) {$0$};
    \node at (0.5,1.5) {$0$}; \node at (1.5,1.5) {$w_{2,2}$}; \node at (2.5,1.5) {$w_{2,1}$};
    \node at (0.5,0.5) {$0$}; \node at (1.5,0.5) {$0$}; \node at (2.5,0.5) {$w_{2,2}$};
    % F1 cinza (linha 3: w_{1,1}, w_{1,2}, 0, 0)
    \node[anchor=south, font=\scriptsize] at (5.5,4.2) {$\mathbf{F}_1$};
    \fill[black!15] (4,0) rectangle (7,4);
    \draw[step=1, gray!60] (4,0) grid (7,4);
    \node at (4.5,3.5) {$w_{1,1}$}; \node at (5.5,3.5) {$0$}; \node at (6.5,3.5) {$0$};
    \node at (4.5,2.5) {$w_{1,2}$}; \node at (5.5,2.5) {$w_{1,1}$}; \node at (6.5,2.5) {$0$};
    \node at (4.5,1.5) {$0$}; \node at (5.5,1.5) {$w_{1,2}$}; \node at (6.5,1.5) {$w_{1,1}$};
    \node at (4.5,0.5) {$0$}; \node at (5.5,0.5) {$0$}; \node at (6.5,0.5) {$w_{1,2}$};
    % F2 laranja (linha 2: zeros)
    \node[anchor=south, font=\scriptsize] at (9.5,4.2) {$\mathbf{F}_2$};
    \fill[orange!30] (8,0) rectangle (11,4);
    \draw[step=1, gray!60] (8,0) grid (11,4);
    \foreach \r in {0.5,1.5,2.5,3.5} \foreach \c in {8.5,9.5,10.5} \node at (\c,\r) {$0$};
    % F3 verde (linha 1: zeros)
    \node[anchor=south, font=\scriptsize] at (13.5,4.2) {$\mathbf{F}_3$};
    \fill[green!25] (12,0) rectangle (15,4);
    \draw[step=1, gray!60] (12,0) grid (15,4);
    \foreach \r in {0.5,1.5,2.5,3.5} \foreach \c in {12.5,13.5,14.5} \node at (\c,\r) {$0$};
\end{tikzpicture}
\vspace{0.3em}
\begin{itemize}
    \item $\mathbf{F}_0$ vem da última linha do filtro paddado, $\mathbf{F}_1$ da penúltima, e assim por diante.
\end{itemize}
\end{frame}

\begin{frame}{\texttt{im2col}: Passo 5 -- Matriz de bloco Toeplitz}
\begin{itemize}
    \item Empilhamos as $F_k$ em uma matriz por blocos. O tamanho final é $16 \times 9$.
\end{itemize}
\vspace{0.4em}
\begin{columns}[T]
    \begin{column}{.45\linewidth}
        \centering
        \begin{equation*}
        \text{Bloco Toeplitz} =
        \begin{bmatrix}
        \textcolor{PastelBlue2!60!black}{\mathbf{F}_0} & 0 & 0 \\
        \textcolor{black!50}{\mathbf{F}_1} & \textcolor{PastelBlue2!60!black}{\mathbf{F}_0} & 0 \\
        \textcolor{orange!70!black}{\mathbf{F}_2} & \textcolor{black!50}{\mathbf{F}_1} & \textcolor{PastelBlue2!60!black}{\mathbf{F}_0} \\
        \textcolor{green!50!black}{\mathbf{F}_3} & \textcolor{orange!70!black}{\mathbf{F}_2} & \textcolor{black!50}{\mathbf{F}_1}
        \end{bmatrix}
        \end{equation*}
    \end{column}
    \begin{column}{.55\linewidth}
        \centering
        \begin{tikzpicture}[scale=0.28, every node/.style={font=\tiny}]
            % 16x9 grid: 4 block rows x 3 block cols, each block 4x3
            % Block (br, bc) at position (bc*3, (3-br)*4)
            % br=0 -> top, br=3 -> bottom
            % Top row: F_0, 0, 0
            \fill[PastelBlue2!35] (0,12) rectangle (3,16);
            \fill[gray!10] (3,12) rectangle (6,16);
            \fill[gray!10] (6,12) rectangle (9,16);
            % Second row: F_1, F_0, 0
            \fill[black!15] (0,8) rectangle (3,12);
            \fill[PastelBlue2!35] (3,8) rectangle (6,12);
            \fill[gray!10] (6,8) rectangle (9,12);
            % Third: F_2, F_1, F_0
            \fill[orange!30] (0,4) rectangle (3,8);
            \fill[black!15] (3,4) rectangle (6,8);
            \fill[PastelBlue2!35] (6,4) rectangle (9,8);
            % Bottom: F_3, F_2, F_1
            \fill[green!25] (0,0) rectangle (3,4);
            \fill[orange!30] (3,0) rectangle (6,4);
            \fill[black!15] (6,0) rectangle (9,4);
            \draw[step=1, gray!50] (0,0) grid (9,16);
            \draw[very thick, gray!80] (0,0) rectangle (9,16);
            \draw[very thick, gray!80] (3,0) -- (3,16);
            \draw[very thick, gray!80] (6,0) -- (6,16);
            \draw[very thick, gray!80] (0,4) -- (9,4);
            \draw[very thick, gray!80] (0,8) -- (9,8);
            \draw[very thick, gray!80] (0,12) -- (9,12);
            \node[anchor=east] at (-0.3,8) {\scriptsize $16$};
            \node[anchor=north] at (4.5,-0.3) {\scriptsize $9$};
        \end{tikzpicture}
    \end{column}
\end{columns}
\vspace{0.3em}
\begin{itemize}
    \item Colunas $=$ linhas da imagem original ($3$); linhas $=$ linhas da imagem de saída vezes a largura ($4 \cdot 4 = 16$).
\end{itemize}
\end{frame}

\begin{frame}{\texttt{im2col}: Passo 6 -- Vetorização da entrada}
\begin{itemize}
    \item Vetorizamos $\mathbf{X}$ concatenando as linhas \textbf{de baixo para cima}:
\end{itemize}
\vspace{0.4em}
\centering
\begin{tikzpicture}[every node/.style={font=\footnotesize}, scale=0.7]
    % X 3x3 azul
    \node at (-0.7,1.5) {$\mathbf{X} = $};
    \fill[PastelBlue2!35] (0,0) rectangle (3,3);
    \draw[step=1, gray!60] (0,0) grid (3,3);
    \node at (0.5,2.5) {$x_{1,1}$}; \node at (1.5,2.5) {$x_{1,2}$}; \node at (2.5,2.5) {$x_{1,3}$};
    \node at (0.5,1.5) {$x_{2,1}$}; \node at (1.5,1.5) {$x_{2,2}$}; \node at (2.5,1.5) {$x_{2,3}$};
    \node at (0.5,0.5) {$x_{3,1}$}; \node at (1.5,0.5) {$x_{3,2}$}; \node at (2.5,0.5) {$x_{3,3}$};
    % Arrow
    \draw[->, gray!70, thick] (3.4,1.5) -- (5.1,1.5);
    % Vector 9x1
    \fill[PastelBlue2!35] (5.5,-3.0) rectangle (6.5,6.0);
    \draw[gray!60] (5.5,-3.0) rectangle (6.5,6.0);
    \foreach \y in {-2,-1,0,1,2,3,4,5}{
        \draw[gray!60] (5.5,\y) -- (6.5,\y);
    }
    \node at (6.0,5.5) {$x_{3,1}$};
    \node at (6.0,4.5) {$x_{3,2}$};
    \node at (6.0,3.5) {$x_{3,3}$};
    \node at (6.0,2.5) {$x_{2,1}$};
    \node at (6.0,1.5) {$x_{2,2}$};
    \node at (6.0,0.5) {$x_{2,3}$};
    \node at (6.0,-0.5) {$x_{1,1}$};
    \node at (6.0,-1.5) {$x_{1,2}$};
    \node at (6.0,-2.5) {$x_{1,3}$};
\end{tikzpicture}
\end{frame}

\begin{frame}{\texttt{im2col}: Passo 7 -- Multiplicação e \textit{reshape}}
\begin{itemize}
    \item Multiplicamos a matriz de bloco Toeplitz pelo vetor de entrada e fazemos \textit{reshape} para a forma da saída:
\end{itemize}
\begin{equation*}
    \underbrace{\mathbf{T}}_{16 \times 9} \cdot \underbrace{\mathbf{x}}_{9 \times 1} = \underbrace{\mathbf{y}}_{16 \times 1}
    \;\;\xrightarrow{\;\text{reshape}\;}\;\;
    \underbrace{\mathbf{Y}}_{4 \times 4}
\end{equation*}
\begin{itemize}
    \item Cada linha do produto $\mathbf{T}\mathbf{x}$ já contém o resultado de uma posição da convolução; basta reorganizar em uma matriz $4\times 4$.
    \item O resultado é exatamente igual ao da convolução tradicional, calculado em \textbf{uma única multiplicação de matrizes}.
\end{itemize}
\end{frame}

\begin{frame}{Por que \texttt{im2col} é a escolha prática}
\begin{itemize}
    \item Reduz a convolução a uma multiplicação de matrizes (\texttt{GEMM}), o tipo de operação que GPUs e bibliotecas BLAS otimizam ao máximo.
    \item Permite usar a mesma rotina para qualquer combinação de \textit{kernel}, \textit{stride} e \textit{padding}, alterando apenas o layout dos dados.
    \item Custa memória extra para armazenar as ``cópias'' das janelas, mas o ganho em velocidade compensa em larga escala.
\end{itemize}
\end{frame}

%========================================================================
\section{Referências}
%========================================================================

\begin{frame}[allowframebreaks]{Referências}
\begin{itemize}
    \item Sugere se \textbf{fortemente} a leitura do Capítulo 10 de \textit{Understanding Deep Learning}\footfullcite{prince2023understanding}.
    \begin{itemize}
        \item Disponível em \url{https://udlbook.github.io/udlbook/}.
    \end{itemize}
    \item Aula baseada também nos seguintes materiais:
    \begin{itemize}
        \item Capítulo 10 de \textit{Understanding Deep Learning}, Simon J. D. Prince.
        \item \textit{Machine Learning} @ Berkeley.
        \item \url{https://github.com/alisaaalehi/convolution_as_multiplication}.
        \item \url{https://github.com/vdumoulin/conv_arithmetic}.
    \end{itemize}
\end{itemize}
\printbibliography
\end{frame}

\end{document}
